% part: intuitionistic-logic
% chapter: introduction
% section: axiomatic-derivations

\documentclass[../../../include/open-logic-section]{subfiles}

\begin{document}

\olfileid{int}{int}{axd}

\olsection{స్వీకృతాధారిత \usetoken{వ్యుత్పత్తులు}{derivation}}

అంతఃప్రజ్ఞావాద ప్రతిజ్ఞావాక్య తర్కానికి
స్వీకృతాధారిత \tetoken{వ్యుత్పత్తులు}{!!{derivation}s}
భావనాపరంగా అతి సరళమైనవి; చారిత్రకంగా తొలి
\tetoken{వ్యుత్పత్తి}{!!{derivation}} వ్యవస్థలూ ఇవే.
ఇవి సాంప్రదాయిక ప్రతిజ్ఞావాక్య తర్కంలోలాగే
పనిచేస్తాయి.

\begin{defn}[\tetoken{వ్యుత్పాద్యత}{!!^{derivability}}]
$\Gamma$ అనేది $\Lang L$ యొక్క
\tetoken{సూత్రాల}{!!{formula}s} సమితి అయితే,
$\Gamma$ నుంచి ఒక \emph{\tetoken{వ్యుత్పత్తి}{!!{derivation}}}
అనేది \tetoken{సూత్రాల}{!!{formula}s}
పరిమిత క్రమం $!A_1$, \dots,~$!A_n$.
ప్రతి $i \le n$కు కింది వాటిలో ఒకటి జరుగుతుంది:
\begin{enumerate}
\item $!A_i \in \Gamma$; లేదా
\item $!A_i$ ఒక స్వీకృతం; లేదా
\item $j < i$, $k <
  i$ అయ్యే ఏవో $!A_j$, $!A_k$ నుంచి మోడస్
  పోనెన్స్ ద్వారా $!A_i$ వస్తుంది; అంటే
  $!A_k \ident !A_j \lif !A_i$.
\end{enumerate}
\end{defn}

\begin{defn}[స్వీకృతాలు]
అంతఃప్రజ్ఞావాద ప్రతిజ్ఞావాక్య తర్కపు
\emph{స్వీకృతాల} సమితి~$\PAx$, కింది రూపాల్లోని
అన్ని \tetoken{సూత్రాలతో}{!!{formula}s} కూడి ఉంటుంది:
\begin{align}
  & (!A \land !B) \lif !A \ollabel{ax:land1}\\
  & (!A \land !B) \lif !B \ollabel{ax:land2}\\
  & !A \lif (!B \lif (!A \land !B)) \ollabel{ax:land3}\\
  & !A \lif (!A \lor !B) \ollabel{ax:lor1}\\
  & !A \lif (!B \lor !A) \ollabel{ax:lor2}\\
  & (!A \lif !C) \lif ((!B \lif !C) \lif ((!A \lor !B) \lif !C)) \ollabel{ax:lor3}\\
  & !A \lif (!B \lif !A) \ollabel{ax:lif1}\\
  & (!A \lif (!B \lif !C)) \lif ((!A \lif !B) \lif (!A \lif !C)) \ollabel{ax:lif2}\\
  & \lfalse \lif !A \ollabel{ax:lfalse1}
\end{align}
\end{defn}

\begin{defn}[\tetoken{వ్యుత్పాద్యత}{!!^{derivability}}]
\tetoken{ఒక సూత్రం}{!!^a{formula}}~$!A$కు,
$\Gamma$ నుంచి మొదలై $!A$తో ముగిసే
\tetoken{ఒక వ్యుత్పత్తి}{!!a{derivation}} ఉంటే,
దాన్ని $\Gamma$ నుంచి \emph{\tetoken{వ్యుత్పాద్యం}{!!{derivable}}}
అంటాం; $\Gamma \Proves !A$ అని వ్రాస్తాం.
\end{defn}

\begin{defn}[సిద్ధాంతాలు]
\tetoken{ఒక సూత్రం}{!!^a{formula}}~$!A$కు
శూన్య సమితి నుంచి~$!A$కు \tetoken{ఒక వ్యుత్పత్తి}{!!a{derivation}}
ఉంటే, అది ఒక \emph{సిద్ధాంతం}. $!A$ సిద్ధాంతమైతే
$\Proves !A$, కాకపోతే $\Proves/ !A$ అని వ్రాస్తాం.
\end{defn}

\begin{prop}
  అంతఃప్రజ్ఞావాద తర్కపు స్వీకృతాధారిత
  \tetoken{వ్యుత్పత్తి}{!!{derivation}} వ్యవస్థలో
  $\Gamma \Proves !A$ అయితే, సాంప్రదాయిక తర్కంలోనూ
  $\Gamma \Proves !A$. ముఖ్యంగా $!A$ ఒక
  అంతఃప్రజ్ఞావాద సిద్ధాంతమైతే, అది సాంప్రదాయిక
  సిద్ధాంతం కూడా.
\end{prop}

\begin{proof}
  ప్రతి అంతఃప్రజ్ఞావాద స్వీకృతం సాంప్రదాయిక
  స్వీకృతం కూడా. కాబట్టి అంతఃప్రజ్ఞావాద
  తర్కంలోని ప్రతి \tetoken{వ్యుత్పత్తి}{!!{derivation}}
  సాంప్రదాయిక తర్కంలోనూ
  \tetoken{ఒక వ్యుత్పత్తి}{!!a{derivation}}.
\end{proof}

\end{document}
